% =====================================================================
%  LAMPIRAN A — DATA PENDUKUNG
% =====================================================================

\chapter{Data Pendukung}
\label{lam:a}

\section{Contoh}
\label{lam:a.1}

Tabel~\ref{tab:lam.a1} dan Persamaan~\eqref{eq:lam.a1}.

\begin{thesistable}
\begin{longtable}{@{}l r@{}}
  \caption{Data pendukung}
  \label{tab:lam.a1}\\
  \toprule
  \textbf{Pengamatan} & \textbf{Nilai} \\
  \midrule
  \endfirsthead
  \multicolumn{2}{@{}l}{\footnotesize\itshape Lanjutan Tabel~\thetable}\\
  \toprule
  \textbf{Pengamatan} & \textbf{Nilai} \\
  \midrule
  \endhead
  \bottomrule
  \endfoot
  \bottomrule
  \endlastfoot
  [Pengamatan 1] & [X,XX] \\
  [Pengamatan 2] & [X,XX] \\
\end{longtable}
\end{thesistable}

\begin{equation}
  \bar{x} = \frac{1}{n}\sum_{i=1}^{n} x_i
  \label{eq:lam.a1}
\end{equation}
